% id: 204-07
% book: 베이직쎈 대수 · 12 수학적 귀납법
% section: 실전 감각 UP
% figure: none
% answer: 
% answer_source: 
% variant_level: 0 (원본 전사)
% 이 파일은 build.mjs 가 items.json 에서 만든다. 고칠 때는 items.json 을 고친다.
\smash{\raisebox{1.5mm}{\dmkichul{베이직쎈 대수 204쪽 07번}}}
홀수인 자연수 $n$에 대한 명제 $p(n)$이 참임을 수학적 귀납법을 이용하여 증명하려고 한다. 반드시 증명해야 할 것을 \textbf{보기}에서 있는 대로 고른 것은?\bogi{ㄱ. $p(1)$이 참이다.\\ ㄴ. $p(2)$가 참이다.\\ ㄷ. $p(k)$가 참이면 $p(k+1)$이 참이다.\\ ㄹ. $p(k)$가 참이면 $p(k+2)$가 참이다.}
\dmchoicesauto{}{ㄱ, ㄷ}{ㄱ, ㄹ}{ㄴ, ㄷ}{ㄱ, ㄴ, ㄹ}{ㄴ, ㄷ, ㄹ}
